% SOURCES: external (you must verify claims first-hand before submission);
%          draft/Features/*.md for how Draw-to-Web positions itself.
% GUIDANCE: The marking scheme wants a review of SIMILAR projects with the
% advantages and shortcomings of EACH, then a positioning of your project.
% Everything below is a SCAFFOLD with candidate points — VERIFY each claim
% against the current product before you cite it; these tools change often.

\chapter{Reference Study}
\label{ch:refstudy}

\section{Survey of Similar Projects}
% DRAFT SCAFFOLD (verify every claim):

\subsection{Webflow}
A browser-based visual builder that emphasises design control and clean output
\cite{webflow}.
\textbf{Advantages (verify):} strong CSS-level control; publishes reasonably
semantic markup; large template ecosystem.
\textbf{Shortcomings (verify):} subscription/hosting lock-in; export can be
gated behind paid tiers; not a local desktop tool; no accessibility hard gate.

\subsection{Framer}
A design-first site builder with rich interactions \cite{framer}.
\textbf{Advantages (verify):} excellent motion/prototyping; fast to a polished
result.
\textbf{Shortcomings (verify):} output leans on a proprietary runtime; less
control over raw semantic HTML; hosting-oriented.

\subsection{Wix}
A mass-market website builder \cite{wix}.
\textbf{Advantages (verify):} very low barrier to entry; all-in-one hosting.
\textbf{Shortcomings (verify):} historically heavy, non-portable output;
limited semantic/accessibility guarantees; strong platform lock-in.

\subsection{Elementor (WordPress)}
A WordPress page builder \cite{elementor}.
\textbf{Advantages (verify):} huge ecosystem; visual editing over WordPress.
\textbf{Shortcomings (verify):} requires a WordPress stack; markup/DOM weight;
plugin-dependent.

\subsection{Bootstrap Studio / Pinegrow}
Desktop visual editors that emit real HTML/CSS \cite{bootstrapstudio,pinegrow}.
\textbf{Advantages (verify):} local, code-visible, export-friendly.
\textbf{Shortcomings (verify):} framework-coupled (Bootstrap); steeper learning
curve; no built-in accessibility export gate; no tokens-first theming.

\section{Comparative Analysis}
% DRAFT SCAFFOLD: fill each cell after you verify. Keep the dimensions that
% distinguish Draw-to-Web (local, semantic, a11y-gated, tokens-first, portable).
\begin{table}[H]
  \centering
  \caption{Comparison of visual web builders (VERIFY before submission).}
  \label{tab:comparison}
  \small
  \begin{tabularx}{\linewidth}{@{}l c c c c c@{}}
    \toprule
    \textbf{Dimension} & \textbf{Draw-to-Web} & \textbf{Webflow} & \textbf{Framer} & \textbf{Wix} & \textbf{Elementor} \\
    \midrule
    Local desktop app        & Yes & No  & No  & No  & No  \\
    Semantic HTML output      & Yes & \textasciitilde & \textasciitilde & No  & \textasciitilde \\
    Accessibility hard gate   & Yes & No  & No  & No  & No  \\
    Tokens-first theming      & Yes & \textasciitilde & \textasciitilde & No  & No  \\
    JS-free output option     & Yes & No  & No  & No  & No  \\
    Portable (no lock-in)     & Yes & \textasciitilde & No  & No  & No  \\
    Headless/scriptable (MCP) & Yes & No  & No  & No  & No  \\
    \bottomrule
  \end{tabularx}
\end{table}
% TODO: replace Yes/No/~ with verified findings; add a paragraph on how the gaps
% above motivate Draw-to-Web's design.

\section{Positioning of Draw-to-Web}
% DRAFT (review):
Against this landscape, Draw-to-Web's distinguishing bets are: a \emph{local}
desktop tool that treats output quality as a first-class requirement (semantic,
deterministic, accessibility-gated), a \emph{tokens-first} styling model, an
\emph{opt-in} runtime (down to JavaScript-free output), and a \emph{headless}
MCP interface to the same pipeline. % TODO: tie back to the shortcomings above.
